\documentclass[sigconf]{acmart}

%% EVERY number in this paper - prose, tables, and abstract alike - is a LaTeX macro defined in
%% figures/variables.tex, generated by research/analysis/paper_variables.py. Nothing below should
%% contain a bare numeric literal that came from a measurement: a number appearing in both a table
%% and the prose discussing it used to be typed twice, so a data refresh could update one and miss
%% the other. See that script's module doc comment for which blocks are read back from a
%% measurement artifact and which are transcribed with their provenance recorded alongside them.
%%
%% The empirical-study numbers (Table 1, corpus size, bytes-AST correlation) were measured over
%% the full \NumRepos{}-repository list on 2026-09-07; see research/data/corpus_stats/PROVENANCE.md
%% for the run. They are no longer the 100-repository `small` sample, and this comment said
%% otherwise until 2026-09-09.
% Auto-generated by research/analysis/file_stats.py. Do not edit by hand - regenerate
% instead: uv run ./analysis/file_stats.py <stats.sqlite> (from research/), then copy
% into research/papers/introductory-paper/figures/ (see that paper's own README.md, or
% make introductory-paper-empirical).
\newcommand{\NumRepos}{100}
\newcommand{\NumFiles}{1{,}347{,}490}
\newcommand{\NumFilesMillions}{1.35}
\newcommand{\NumLanguages}{30}
\newcommand{\CorrelationR}{0.8794}
\newcommand{\BytesPFifty}{1{,}508}
\newcommand{\BytesPNinety}{13{,}945}
\newcommand{\BytesPNinetyNine}{90{,}132}
\newcommand{\BytesPNineNineNine}{464{,}325}
\newcommand{\BytesMax}{23{,}949{,}786}
\newcommand{\LocPFifty}{42}
\newcommand{\LocPNinety}{378}
\newcommand{\LocPNinetyNine}{2{,}244}
\newcommand{\LocPNineNineNine}{7{,}818}
\newcommand{\LocMax}{65{,}563}
\newcommand{\AstPFifty}{224}
\newcommand{\AstPNinety}{2{,}868}
\newcommand{\AstPNinetyNine}{18{,}553}
\newcommand{\AstPNineNineNine}{70{,}828}
\newcommand{\AstMax}{573{,}334}


%% Result boxes. Each research question stated in the introduction is answered by a correspondingly
%% numbered research answer, RA1.1--RA4.2, set in one of these at the point where its supporting
%% measurement is reported. A box holds the answer statement only - two or three sentences, every
%% measured number a macro from figures/variables.tex - with the argument behind it following as
%% ordinary prose. Keeping the boxes short is a two-column-layout constraint, not a style
%% preference: a boxed half-page of prose splits badly across a column break.
\usepackage[most]{tcolorbox}
\usepackage{array}
\newtcolorbox{resultbox}[1]{%
  breakable,
  enhanced,
  colback=black!3,
  colframe=black!55,
  colbacktitle=black!10,
  coltitle=black,
  fonttitle=\bfseries,
  boxrule=0.4pt,
  left=5pt, right=5pt, top=4pt, bottom=4pt,
  title={#1},
}

\setcopyright{acmlicensed}
\copyrightyear{2026}
\acmYear{2026}
\acmDOI{XXXXXXX.XXXXXXX}
\acmConference[Conference'26]{Make sure to enter the correct conference title from your rights
  confirmation email}{June 03--05, 2026}{Ludbreg, Croatia}
\acmISBN{978-1-4503-XXXX-X/26/06}

\begin{CCSXML}
<ccs2012>
   <concept>
       <concept_id>10011007.10011006.10011073</concept_id>
       <concept_desc>Software and its engineering~Software maintenance tools</concept_desc>
       <concept_significance>500</concept_significance>
       </concept>
   <concept>
       <concept_id>10011007.10011074.10011099.10011693</concept_id>
       <concept_desc>Software and its engineering~Empirical software validation</concept_desc>
       <concept_significance>300</concept_significance>
       </concept>
 </ccs2012>
\end{CCSXML}

\ccsdesc[500]{Software and its engineering~Software maintenance tools}
\ccsdesc[300]{Software and its engineering~Empirical software validation}

\keywords{source code differencing, syntax-aware diffing, empirical software engineering, source code tree edit distance}

\title{CodeDiff: A Fast, Robust, Syntax-Aware Code Diffing Tool}

\author{Marko Ivankovi\'{c}}
\affiliation{%
  \institution{Universit\"{a}t Passau}
  \city{Passau}
  \country{Germany}
}
\email{marko@ivankovic.me}

\author{Gordon Fraser}
\affiliation{%
  \institution{Universit\"{a}t Passau}
  \city{Passau}
  \country{Germany}
}
\email{gordon.fraser@uni-passau.de}

\begin{document}

\begin{abstract}
Code diffing, the process of comparing two pieces of source code and computing the set of
operations that transforms one into the other, underlies modern version control and, especially,
modern code review. During a review, we read the diff to see what changed and what stayed the
same, so the quality of the diff directly shapes the ease and quality of the review. A
syntax-aware diff maps one abstract syntax tree (AST) onto another. AST diffing algorithms assume the following: that the best
diff can be expressed with the four standard operations (insert, delete, update, move) and is
the cheapest one under a unit cost model with a free move, that computing it is possible within
an interactive wall-clock budget, and that the AST node mapping can be
visualised without losing quality. We test each assumption on a corpus of real code changes, to measure the viability of AST
diffing.

We built a corpus of \NumFixtures{} real-world changes in \NumFixtureLanguages{} languages, each
carrying a human-authored AST mapping and \PaintingFixtures{} of them a human-painted rendering.
On \AmbiguityListPct{}\% of changes the best possible correspondence requires a mapping that the
common operations (insert, delete, update and move) cannot express; on at least \CostHumanHigherPct{}\%
the human-preferred mapping is strictly costlier than the cheapest AST diff; and
\PaintingDualPct{}\% of painted changes admit at least two equally correct text renderings of one
settled node mapping. \InvisibleNodeShare{}\% of AST nodes carry no text of their own and never reach
the screen, so a node-level metric gives a third of its weight where no reader can see. A single
whole-tree tree-edit-distance computation completes within one second for only \RqOneCodePct{}\%
of \RqOnePairsAttempted{} stratified-sampled code file pairs.
Because no two tools report a diff in the same form, we compare them on the share of changes
each renders perfectly, marking exactly the lines the human mapping marks. Against this corpus,
no configuration of seven established tools renders more than \GumTreePerfectPct{}\% of changes
perfectly, and plain line-based \texttt{diff} renders \UnixDiffPerfectPct{}\%.

Finally, as a tool contribution, we present \textsc{CodeDiff}, which combines existing
state-of-the-art algorithms. It maps \NodeAccuracyPct{}\% of AST nodes correctly and, on the
same basis as those ten configurations, renders \CodeDiffPerfectPct{}\% of changes perfectly. It
completes \RobustnessFullCompletedPct{}\% of the \RobustnessFullAttempted{} file modifications
in the corpus's recent history without a panic.
\end{abstract}

\maketitle

\section{Introduction}
\label{sec:intro}

\begin{figure}[t]
  \small
  \begin{tabular}{@{}p{0.45\linewidth}@{\hspace{0.05\linewidth}}p{0.45\linewidth}@{}}
    \textsc{Before} & \textsc{After} \\
    \midrule
    \begin{minipage}[t]{\linewidth}
\begin{verbatim}
def save(rec):
  db.begin()
  db.put(rec)
  db.commit()
\end{verbatim}
    \end{minipage}
    &
    \begin{minipage}[t]{\linewidth}
\begin{verbatim}
def save(rec):
  if rec.dirty:
    db.begin()
    db.put(rec)
    db.commit()
\end{verbatim}
    \end{minipage} \\
  \end{tabular}
  \caption{Wrapping a block in a condition. Only whitespace (indentation) changed, but a
  line-based diff reports all of it as deleted and reinserted. An AST diff can report one
  insertion and a move.}
  \label{fig:motivation}
\end{figure}

Code diffing, the process of comparing two pieces of source code and computing the set of
operations that transforms one into the other, underlies modern version control and, especially,
modern code review. During a review, we read the diff to see what changed and what stayed
the same, so the quality of the diff directly shapes the ease and quality of the review.

The standard diff used by almost all developers and platforms today is line-based: the file is
treated as a sequence of text lines, and the tool computes a shortest edit script between the two
sequences using three operations, insert, delete, and update. This approach is computationally
cheap and fast even on modest hardware, but most implementations offer no move operation, and the
tools that do support move use it only in limited ways. A whole class of common
changes—extracting code into a reusable function, wrapping a block in a condition, reordering
declarations—is a move of existing code, and without a move operation every such
change must be rendered as a full deletion plus a full insertion. Reconstructing the
correspondence between the deleted and inserted halves is then left to the reviewer, by eye,
which is the kind of mechanical, attention-taxing work humans make mistakes when doing.
Figure~\ref{fig:motivation} shows a small example of the problem: wrapping three unchanged
lines in a conditional, which a line-based tool reports as six lines deleted and seven
inserted, and which a reader has to diff again in their head to see that only the wrapper is new.

One way to compute diffs with a genuine move operation is to compare the code's abstract syntax
tree (AST) instead of its raw text. An AST diff maps nodes of one tree onto nodes of the other; a
move is a mapping whose nodes match but whose ancestry changed. Tree diffing is a well-studied
field, and Section~\ref{sec:background} surveys the work this paper builds on. In practice,
however, tree diffing has seen little adoption, and two obstacles stand between the algorithms and
a tool developers use. The first is runtime: pure exact tree-diffing algorithms are $O(n^3)$ in the number of nodes. The second is that an AST diff is not directly usable by a developer. It has
to be visualised, usually overlaid over the textual representation of the source code, and for any
given AST diff there are several possible visualisations---so a tool must choose one, and the quality of
that choice determines adoption as much as the mapping underneath it does.

What makes a diff good is what a reviewer does with it. A reviewer reads the diff to decide what
changed and whether the change is right, and every place where the diff misstates the
correspondence between the two versions is a place where the reviewer has to recover that
correspondence by hand, or fails to. Four properties follow, and each is a question about real
code changes rather than about any one tool.

\begin{itemize}
  \item \textbf{Expressible.} The diff must be able to say what happened. If the true
    correspondence between the two versions cannot be written in the operations a tool has, no
    tool of that kind can show it, however well it searches. We ask how often that is the case
    (RQ1.1).
  \item \textbf{Agreed.} Where several correspondences are possible, the diff must pick the one a
    reader would, and the reader's choice is not always the cheapest. We ask how often the human
    choice is costlier than the optimum (RQ1.2) and how often there is no single correct choice at
    all (RQ3.1).
  \item \textbf{Timely.} The diff must arrive before the reader's attention lapses. We ask how much
    of a real corpus an exact whole-tree computation can serve within one
    second\footnote{One second is where Nielsen's response-time limits put the boundary between an
    interaction that keeps a user's train of thought and one that visibly interrupts
    it~\cite{nielsen-usability}.} (RQ2), and how fast the tools that exist are (RQ4.2).
  \item \textbf{Faithful on screen.} The reader sees text, not nodes. What the rendering marks has
    to be what the mapping says, and a mapping is only as good as the part of it that reaches the
    screen. We ask how often one settled mapping admits more than one correct rendering (RQ3.2)
    and how much of the AST carries visible text at all (RQ3.3).
\end{itemize}

The last question is whether the tools that exist today have these properties: how accurately they
find the mapping a human would (RQ4.1) and how fast they are (RQ4.2). We do not assume the answer.
In particular, we do not take it as given that a syntax-aware diff serves a reader better than a
line-based one. That is what RQ4.1 measures, with line-based \texttt{diff} scored on the same basis
as every AST-aware tool, and Section~\ref{sec:tools} reports where each kind does better.

%% Each research question is stated in full in the section that answers it, next to its result
%% box. The introduction names them by number under the property each one tests; keep the two in
%% step when a question is renumbered or reworded.
Each question is answered in its own section: Section~\ref{sec:viability} takes RQ1,
Section~\ref{sec:speed} RQ2, Section~\ref{sec:uniqueness} RQ3 and Section~\ref{sec:tools} RQ4.
Section~\ref{sec:codediff} then reads every answer as a design decision in a tool.

This paper makes the following contributions:

\begin{itemize}
  \item \textbf{Dataset:} A set of \NumFixtures{} real-world code changes. The changes cover
    \NumFixtureLanguages{} languages and were selected using 3 different selection algorithms from
    \NumRepos{} open source repositories.
    The repositories themselves are a combination of a curated list and every repository used by the Gentoo
    Linux distribution. Every change was solved by hand and the optimal map between the two ASTs is
    available, or if the optimal map is impossible using the common AST diff operations, a best-possible
    solution with a comment explaining the missing part. Additionally, \PaintingFixtures{} of those
    solutions also have the optimal visualisations: the spans of text a developer should see marked,
    independently of which AST nodes carry them. Often, two visualisations are provided: the minimal one
    and the maximal one, covering human preferences between equally valid solutions.
    The dataset, the tool used to create the dataset and the entire history of the dataset is available
    freely for future research. To allow comparison with existing research, the dataset has a
    fourth part: the Defects4J~\cite{defects4j} changes of the AST node-mapping oracle of
    Alikhanifard and Tsantalis~\cite{astdiffbenchmark}, solved by us by hand for comparison with an
    independently authored reference. The dataset is described in detail in
    Section~\ref{sec:dataset}.
  \item \textbf{Evaluation:} Answers to all research questions. An analysis of the real-world code change
    data shape, irrespective of any concrete implementation, and a quantified comparison of several
    state-of-the-art code diffing tools on the human annotated dataset.
  \item \textbf{Tool:} \textsc{CodeDiff}, a new state-of-the-art tool. Combining all current research
    with strict engineering principles and carefully implemented to maximise performance.
\end{itemize}

\section{Background}
\label{sec:background}

This section lists two types of existing work: pure algorithms and implemented tools.

\subsection{Algorithms}

\textbf{Zhang-Shasha}~\cite{zhang-shasha-1989} gives the classic dynamic-programming algorithm for
ordered tree edit distance. Given two rooted, ordered trees, it computes the minimum-cost sequence
of node insertions, deletions, and relabelings that transforms one tree into the other. It is
exact, but its keyroot decomposition does repeated work across overlapping subproblems.

\textbf{APTED}~\cite{apted-2015} is the state-of-the-art exact tree-edit-distance algorithm. It
improves on Zhang-Shasha by choosing, for each subproblem, the cheapest of three decomposition
strategies (left path, right path, or heaviest inner path), rather than a single fixed strategy.
This choice provably minimises the total work across the whole computation. APTED still costs
$O(n^3)$ in the worst case for general trees, the same asymptotic bound as Zhang-Shasha, but its
real-world runtime is markedly lower.

\textbf{Myers' O(ND) algorithm}~\cite{myers-1986} solves a different, more restricted problem: the
shortest edit script between two flat sequences, not two trees. It underlies most line-based diff
tools, including Unix \texttt{diff}. Its cost scales with the size of the edit script, not the size
of the inputs, so it stays fast when two sequences are similar.

\subsection{Tools}

\textbf{Unix \texttt{diff}} is the golden standard. It treats each file as a sequence of lines, computes a shortest edit script
between the two sequences with Myers' algorithm, and reports it as hunks of deleted and inserted
whole lines. It has no move operation and no sub-line detail.

\textbf{\texttt{git diff}} defaults to the \texttt{xdiff} engine, which offers four
algorithms: Myers, the default; \emph{minimal}, which spends more time to find a shortest script; \emph{patience}, which
anchors on lines unique to both sides before recursing; and \emph{histogram}, a faster
generalisation of patience. All four emit the same output form as Unix \texttt{diff}.

\textbf{Neovim's diff mode} (\texttt{nvim -d}) runs the same \texttt{xdiff} line pass as
\texttt{git} and then a second, character-level pass that highlights what changed inside
each paired line.

\textbf{BDiff}~\cite{bdiff} is a text-based tool that reasons about blocks of lines rather
than single lines and reports character ranges within changed lines.

\textbf{GumTree}~\cite{gumtreediff-2014} is a source code AST diff. Instead of one global
tree-edit-distance computation, it runs a top-down phase that greedily matches large
isomorphic subtrees, then a bottom-up phase that matches remaining container nodes by a Dice
coefficient of their already-matched descendants, and a recovery pass that pairs up deleted and
inserted identical subtrees the top-down and bottom-up phases missed. For performance, it also introduces InsertTree and DeleteTree operations
that are equivalent to repeatedly calling Insert or Delete operations on all children in the given
subtree.

\textbf{difftastic}~\cite{difftastic} and \textbf{diffsitter}~\cite{diffsitter} are a newer
generation of tree-sitter-based diffing tools. Both prioritise a diff a human reads
comfortably over an exact node-to-node mapping: difftastic reduces the parse tree to an s-expression
and aligns it with its own structural algorithm, and diffsitter applies a related tree-sitter-based
alignment over a smaller, hand-picked set of compiled-in grammars.

\section{Empirical Dataset}
\label{sec:dataset}

To understand the real-world environment code diffing tools are used in, we built a dataset of \NumFiles{}
files in \NumLanguages{} languages. Of these, we sampled \NumFixtures{} real open source code
changes and examined them in depth.

To ensure good coverage, we built two sets of open-source repositories:

\begin{itemize}
  \item The Curated list - 100 open source repositories, selected manually from several lists of
    most popular repositories, enhanced further by repositories we were already aware of. This biased the list towards active, widely used codebases, with good coverage in
    many programming languages. The Curated list was also used as a smaller list to repeatedly
    test the paper's own codebase, including in automated test infrastructure.
  \item The Full List - \NumRepos{} repositories, created by expanding the Curated list
    with all repositories present in the Gentoo Linux distribution's package list. This includes
    a wide array of codebases of different activity, code quality, userbase and programming language
    mixes.
\end{itemize}

We cloned each repository to a depth of \CorpusCloneDepth{} commits from every main branch. While we cannot guarantee that all repositories
remain accessible, the entire list of repositories, the cloning script and all analysis is available from the paper's repository and can be replicated by readers, if desired.

We fetched, parsed and measured the corpus on a single machine: an \MachineCpu{} with \MachineRamGb{}\,GB of RAM, running \MachineOs{}. The checkouts and the measurement database sit on \MachineDiskCount{} \MachineDiskSizeTb{}\,TB \MachineDiskModel{} drives in \MachineDiskArray{}. We record this because some results are hardware-bound, notably the wall-clock budget we hold whole-tree APTED to in Section~\ref{sec:speed} and the per-tool wall-clock timing of Section~\ref{sec:tools}.

This section describes the overall shape of the resulting dataset.

\subsection{Shape of real-world files}

We used a simple extension- and content-based method to categorise the files. The first thing to note about the dataset is that barely half of all files are code at all (Figure~\ref{fig:tips}). The rest is configuration, data, documentation, and a substantial category of unknown type; based on limited manual inspection we believe most of the unknown-type files to be data. Note however that any ``code'' diffing tool still must support visualising changes to non-code files as well.

Restricted to code files, the top row of Figure~\ref{fig:corpus-shape} shows the distribution of
bytes, lines of code and AST nodes per file. The median code file has \LocPFifty{} lines and
\AstPFifty{} AST nodes, the 99th percentile \LocPNinetyNine{} lines and \AstPNinetyNine{} nodes,
and the largest \LocMax{} lines and \AstMax{} nodes: three orders of magnitude between the
typical file and the tail a tool still has to handle. Files are parsed whatever their size,
except those the statistics flag as generated from a header comment, which are skipped; the
largest files that reach the parser are generated data that carry no such comment---voice-model
parameters, benchmark inputs, parser tables---and every file of that order is one,
which Section~\ref{sec:codediff} returns to.

\begin{figure}
  \centering
  \includegraphics[width=0.95\linewidth]{figures/tips}
  \caption{Distribution of file types across the corpus.}
  \label{fig:tips}
\end{figure}

\begin{figure*}
  \centering
  \includegraphics[width=\linewidth]{figures/corpus_shape}
  \caption{The shape of the corpus, every distribution drawn the same way: the share of the
  population at or below the value on the axis, with the median, 90th and 99th percentiles marked
  on each curve and the maximum in the corner. Top: size of every code file in the corpus, in
  lines, bytes and AST nodes; the flat start of the nodes curve is the files the statistics
  give no node count, which are counted: those they flag as generated and do not parse, those in
  a language without a grammar, and empty ones. Bottom: size of every code-file modification in the
  corpus's recent history, in lines changed per file, lines changed per commit, and the share of
  the file an edit rewrites. The horizontal axes are logarithmic except the last.}
  \label{fig:corpus-shape}
\end{figure*}

Byte size and AST node count are strongly correlated for files within the 99th percentile of
both, Pearson $r = \CorrelationRTrimmed{}$, which simplifies to $|\text{AST}| \approx \text{bytes} /
5$ for quick estimates. Over every file $r$ is \CorrelationR{}: the largest one percent are the
generated tables above, each with a byte-per-node ratio of its own, and a correlation over a
heavy-tailed population follows its largest points.

\subsection{Shape of real-world file edits}
\label{sec:edit-shape}

Across both repository lists we analysed \EditsCommits{} commits. To stop large repositories,
notably the Linux kernel, from skewing the results, each repository was analysed at most
\CorpusCloneDepth{} commits deep. \EditsCodeSharePct{}\% of the files those commits touched
were code files, giving \EditsCodeFileEdits{} code-file edits, of which
\EditsModifiedSharePct{}\% were modifications; the remainder created or deleted a file and so
present no pair of versions to map between.

To allow for comparison with existing research, \texttt{git diff} was used to compute the number of
changed lines.


The bottom row of Figure~\ref{fig:corpus-shape} shows the sizes of real-world edits. The
median file edit changes \EditsLinesPerFilePFifty{} lines and the median commit
\EditsLinesPerCommitPFifty{}; \EditsFileEditsUnderTenPct{}\% of file edits and
\EditsCommitsUnderTenPct{}\% of commits change ten lines or fewer, while the 99th percentiles are
\EditsLinesPerFilePNinetyNine{} and \EditsLinesPerCommitPNinetyNine{} lines. Our results are
broadly aligned with existing research. Arafat and Riehle report that 47.5\% of all commits touch 10 lines or
fewer~\cite{commit-size-distribution-paper}, with the long tail stretching to tens of thousands of
lines. The median edit rewrites only \EditsChurnPFiftyPct{}\% of the file it touches, and
\EditsChurnUnderFivePct{}\% of edits rewrite under 5\%. The tail is the opposite extreme rather
than a larger fraction of the same shape: by the 99th percentile an edit rewrites the file
entirely (\EditsChurnPNinetyNinePct{}\%). A diffing tool therefore has to be good at two
unrelated jobs---finding a handful of changed lines in an otherwise untouched file, which is the
overwhelming majority of edits, and staying usable when nothing is left to anchor on.

\subsection{Shape of optimal file diffs}
\label{sec:diff-shape}

From both Curated and Full list of repositories, we fully randomly sampled \SampleDiffsPerLanguage{}
diffs per language. That gave \SampleCuratedSampled{} diffs from the Curated list across
\SampleCuratedLanguages{} languages, and \SampleFullSampled{} from the Full list across
\SampleFullLanguages{} languages. The Full list's count exceeds \SampleDiffsPerLanguage{} per
language because some languages were drawn a second time after their first rounds were largely
rejected. Every sampled diff was manually inspected for
sampling errors.
At this stage, \SampleRejectedTotal{} diffs were rejected, mostly because programming language detection detected
the wrong language. In particular TypeScript (\SampleRejectedTypeScript{} rejected) and R (\SampleRejectedR{}) file
extensions were used for other languages or data. This resulted in a dataset of \SampleCuratedPromoted{} diffs from the
Curated and \SampleFullPromoted{} diffs from the Full list. As shown in Section~\ref{sec:edit-shape}, a pure random
sample would be heavily slanted towards small changes. To account for this, the fully random sample was further
augmented by a stratified sample based on the lines of code of the files being changed. The larger of the two sizes was used.
The strata were: 1 to 10 lines, 11 to 30, 31 to 100, 101 to 300, 301 to 1000, 1001 to 3000 and
more than 3000 lines of code. In each stratum, \SampleStratifiedPerCell{} diffs were randomly
sampled per language, and inspected as above.\footnote{\label{fn:in-progress}The corpus and this
paper are being written in parallel. Drawing and solving the stratified sample, and hand-annotating
the Defects4J list below, are both still in progress, and the ambiguity pass over those two has not
been made at all. Every figure in this paper is over the \NumFixtures{} changes solved so far, and
every figure that depends on the ambiguity pass is a lower bound for those two datasets.}

A fourth list further extends the corpus. The AST node-mapping oracle of Alikhanifard and
Tsantalis~\cite{astdiffbenchmark} contains the \OracleUnits{} AST mappings for \OracleCases{}
bug-fixing changes of Defects4J~\cite{defects4j}. The file contents its offsets index into come
from Falleri and Martinez' replication package~\cite{gumtreesimple}. We solve those same
\OracleUnits{} diffs by hand as well, so that our own mappings can be read against an independently
authored reference on a corpus that other AST-differencing evaluations already
use.\textsuperscript{\ref{fn:in-progress}} Section~\ref{sec:oracle} describes that oracle and reads
the two annotations against each other.

Each diff was then solved by hand. Solving a diff means writing down, for every node of the
before-AST, which node of the after-AST it corresponds to, or that it has none. That is the same
object an AST diffing algorithm produces, except that the correspondence is decided by a person
reading the two versions rather than by a cost function. The annotator works in a small terminal
tool, published with the corpus, that shows the two trees side by side and lets them walk both
independently and mark a pair of nodes as matched or a single node as inserted or deleted. The
tool classifies each matched pair itself, as identical, updated, or matched with differing
children, by comparing the nodes' text, and it can pair up byte-identical subtrees in one
keystroke; everything else is marked by hand, starting from an empty mapping rather than from any
tool's proposal. A fixture's solution is that list of entries, one file per change, stored next to
the two versions of the file and a provenance note naming the repository, commit and path they
came from.

The criterion is the correspondence a reviewer of that change would want shown, which is not
always the cheapest one; Section~\ref{sec:viability} measures the gap. Where two correspondences
are equally right, both are recorded rather than one chosen (Section~\ref{sec:uniqueness}), and
where no set of one-to-one node pairs can express the change at all, the closest expressible
solution is recorded with a note saying what it misses. For \PaintingFixtures{} changes the
annotator additionally painted the rendering: the spans of text a reader should see marked, chosen
independently of which nodes carry them. The annotators are the authors, who have each at least
15 years of experience in software development and have reviewed tens of thousands of real-world
code changes in industrial environments.


Please note, diffing multiple files and detecting cross-file code moves is out of scope of this paper.

\subsection{An independently authored oracle}
\label{sec:oracle}

Every accuracy figure in this paper is scored against ground truth this paper's authors wrote. The
AST node-mapping oracle of Alikhanifard and Tsantalis~\cite{astdiffbenchmark} is the answer to the
obvious objection. It records, for \OracleCases{} bug-fixing changes of Defects4J over
\OracleUnits{} Java compilation units, which before-side node each after-side node corresponds to:
\OracleRecords{} node correspondences in all, built by running six tools, inspecting every diff by
hand against six stated criteria, and resolving disagreements by vote. It was authored by different people, for a
different tool, on a different AST, before \textsc{CodeDiff} existed. The file contents its
character offsets index into come from Falleri and Martinez' replication
package~\cite{gumtreesimple}.

The oracle's unit is an Eclipse JDT node and ours is a tree-sitter node, so the two meet only where
a record's character span is one some tree-sitter node also has. \OracleResolutionRate{}\% of
records resolve; the remainder are JDT constructs tree-sitter does not model separately, chiefly
the elements inside a Javadoc comment, which tree-sitter reads as a single token. That share is the
part of the oracle we could not ask, and every rate below is over the part we could.

The \AmbiguityDefectsFourJTotal{} compilation units we have solved so far are reported as a fourth
dataset, Defects4J, beside the Curated, Full and Stratified samples throughout, and are included in
every pooled total. One caveat belongs to that dataset and not to the other three: which units are
solved is how far annotation has reached rather than a draw, so a Defects4J figure is a reading of
those units and not an estimate over the \OracleUnits{}.

Read against our own annotation rather than against any tool, the oracle is the only check in this
paper that does not depend on our annotation discipline at all. Over the \OracleHumanUnits{}
compilation units (\OracleHumanCases{} cases) we have mapped by hand so far, our mapping and the
oracle's agree to \OracleHumanPrecision{}\% / \OracleHumanRecall{}\%:
\OracleHumanDisagreements{} disagreeing node pairs out of \OracleHumanMappings{}, and
\OracleHumanStatementDisagreements{} of \OracleHumanStatementMappings{} at statement level. The
control that makes those figures readable is \textsc{CodeDiff} over the same units, which agrees
with the oracle to \OracleHumanCodeDiffPrecision{}\% / \OracleHumanCodeDiffRecall{}\%, or
\OracleHumanCodeDiffDisagreements{} disagreeing pairs, and to
\OracleHumanCodeDiffStatementPrecision{}\% / \OracleHumanCodeDiffStatementRecall{}\% over
\OracleHumanCodeDiffStatementDisagreements{} at statement level.

The difference between the two readings is mostly precision, and reading it that way is what makes
it informative. The two recover almost the same share of what the oracle records
(\OracleHumanRecall{}\% against \OracleHumanCodeDiffRecall{}\%), so neither is systematically
missing pairs the other finds. What separates them is pairs \emph{claimed} and not
corroborated---\OracleHumanFalsePositives{} for our annotation against
\OracleHumanCodeDiffFalsePositives{} for the tool, and \OracleHumanStatementFalsePositives{}
against \OracleHumanCodeDiffStatementFalsePositives{} at statement level. Two independently
authored accounts of the same changes therefore invent about half as many pairs as the tool does,
and under a tenth as many at statement level, which is the shape one would expect if the residue is
the tool's rather than a difference of annotation opinion.

That reading is bounded by which units are solved. They are not a sample: all \OracleUnits{} were
taken into the corpus at once, and the solved ones are simply how far annotation has reached. They
run much shorter than the rest---a median of \OracleHumanRecordsSolved{} oracle records against
\OracleHumanRecordsRest{}---so this is a statement about the smaller changes in the set, and it may
move in either direction as annotation reaches the large files.

\section{Viability}
\label{sec:viability}

A diff that cannot say what happened is wrong before any search begins, and a diff that says
something the reader disagrees with is wrong however cheap it was. This section tests those two
properties, Expressible and Agreed, on the human solutions alone, before any tool is involved.
In Section~\ref{sec:diff-shape} we note that human solvers were not able to solve all diffs using
the four standard AST diff operations. We used that information to answer the following research question:

\textbf{RQ1.1.} What percentage of real-world code changes have an optimal AST node mapping that
cannot be expressed using common AST diff operations (Insert, Delete, Update, Move)?

\AmbiguityListWith{} of \AmbiguityListScored{} solved diffs (\AmbiguityListPct{}\%) require an
$N{:}M$ mapping: one or more nodes from one side map to one or more nodes from the
other.\footnote{The most common cause of an $N{:}M$ mapping is a refactoring, where two or more
similar blocks of code are moved into a common function. The correct solution is the one that maps
all of them as moved.} That rate is over the \AmbiguityListScored{} changes of the Curated and Full
lists. The two lists agree closely, at \AmbiguityCuratedPct{}\% of \AmbiguityCuratedTotal{}
changes on the Curated list and \AmbiguityFullPct{}\% of \AmbiguityFullTotal{} on the Full list
(Table~\ref{tab:expressibility}). Those two are the datasets whose ambiguity pass is finished, and
the pooled rate RA1.1 states is over them. The table also carries the Stratified sample
(\AmbiguityStratifiedPct{}\%) and the Defects4J changes (\AmbiguityDefectsFourJPct{}\%), and a rate
over all \AmbiguityScored{} changes, but all three are lower bounds: those two passes are still
running, and a change nobody has reviewed for ambiguity is recorded as having none. Read the gap
between \AmbiguityListPct{}\% and \AmbiguityAnyPct{}\% as the distance annotation still has to
travel, not as a difference between the datasets.

\begin{table}
  \caption{Changes whose optimal mapping needs an $N{:}M$ site, per dataset of
  Section~\ref{sec:dataset} and pooled.}
  \label{tab:expressibility}
  \begin{tabular}{lrrr}
    \toprule
    Dataset & $n$ & $N{:}M$ changes & Share \\
    \midrule
    Curated & \AmbiguityCuratedTotal{} & \AmbiguityCuratedWith{} & \AmbiguityCuratedPct{}\% \\
    Full & \AmbiguityFullTotal{} & \AmbiguityFullWith{} & \AmbiguityFullPct{}\% \\
    Stratified\textsuperscript{*} & \AmbiguityStratifiedTotal{} & \AmbiguityStratifiedWith{} &
      \AmbiguityStratifiedPct{}\% \\
    Defects4J\textsuperscript{*} & \AmbiguityDefectsFourJTotal{} & \AmbiguityDefectsFourJWith{} &
      \AmbiguityDefectsFourJPct{}\% \\
    \midrule
    \emph{Curated + Full} & \AmbiguityListScored{} & \AmbiguityListWith{} & \AmbiguityListPct{}\% \\
    \emph{All}\textsuperscript{*} & \AmbiguityScored{} & \AmbiguityAnyFixtures{} &
      \AmbiguityAnyPct{}\% \\
    \bottomrule
  \end{tabular}

  \vspace{2pt}
  \footnotesize\textsuperscript{*}The ambiguity pass over the Stratified and Defects4J changes is
  not finished, so these three figures are lower bounds: a change nobody has reviewed for ambiguity
  is recorded as having none. \emph{Curated + Full} is the rate over the two datasets whose pass is
  complete, and is the figure RA1.1 states.
\end{table}

\begin{resultbox}{RA1.1 (Expressibility)}
  \AmbiguityListPct{}\% of changes (\AmbiguityListWith{} of \AmbiguityListScored{}) have optimal
  solutions that cannot be expressed at all using the common AST diff operations (Insert, Delete,
  Update, Move).
\end{resultbox}

Please note that this result has direct implications for other results in this paper. Because
no tool can possibly solve these diffs fully, we used a charitable interpretation when evaluating
tools and exclude those nodes from the evaluation. In the datasets, the nodes are left unmatched for now.

Tree diffing algorithms find the lowest cost diff, i.e. smallest number of operations needed to
transform one tree into another. However, this doesn't necessarily correspond to the optimal
visualisation of the diff overlayed over the text of the source code.

\textbf{RQ1.2.} What percentage of real-world code changes have solutions that humans prefer
whose cost is higher than the lowest cost AST diff, using the commonly accepted costs of AST
diff operations?

The lowest edit distance was computed using a state of the art implementation. Of the \CostScored{} fixtures with a human mapping, the
two costs are equal \CostTieFixtures{} times. In \CostHumanHigherFixtures{} (\CostHumanHigherPct{}\%) the human solution is strictly costlier. The most common reason for this is that humans treat groups of tokens
as one unit, not based on their syntax or semantical meaning, but rather based on higher-level logic.
A common example is the \texttt{return} keyword; if one function was added and another deleted,
matching the two leads to an optimal tree edit distance, but a human would find the match confusing because they typically
treat the keyword as ``belonging'' to the logic of the function containing it. Existing tools do take
notice of this. GumTree~\cite{gumtreediff-2014} avoids matching subtrees smaller than a
threshold size to avoid such spurious mappings.

\begin{resultbox}{RA1.2 (Cost)}
  On at least \CostHumanHigherPct{}\% of the changes in the corpus (\CostHumanHigherFixtures{} of
  \CostScored{}), the human-authored mapping costs strictly more, under unit costs for insert,
  delete and update and a free move, than a mapping a tree edit distance algorithm found for the
  same change.
\end{resultbox}

\section{Speed}
\label{sec:speed}

A correct diff that arrives late is a diff the reader has stopped waiting for. This section
tests the Timely property on the exact algorithm alone, so that the cost of correctness is known
before any tool's shortcuts are measured against it.
The asymptotic complexity of APTED and Zhang-Shasha both is $O(n^3)$ for general trees, where $n$
is the number of nodes in the tree. However, APTED performs considerably better in real trees
because it is less sensitive to the tree shape. Given this existing direction in the literature,
we wanted to further focus on source-code trees specifically and ask the following question:

\textbf{RQ2.} What percentage of real-world source-code changes can a single, whole-tree
tree-edit-distance computation complete within one second?

One second was chosen as a budget because it is a practical limit that can be computed,
and because it is where interaction stops feeling immediate. Nielsen's response-time
limits~\cite{nielsen-usability} put one second at the boundary at which a user's flow of thought
stays uninterrupted but the delay becomes noticeable; below about 100\,ms a response reads as
instantaneous. From an industry perspective, one second is the high end of acceptable.
Developers have had decades to get used to instant diffs from line-based tools, so any AST diffing tool that expects to be adopted must meet a similar bar.

We sampled \RqOnePairsAttempted{} real (before, after) file pairs from single-parent commits in the
dataset, stratified per language and per size bucket. The buckets used in this stratification were the same
as defined in Section~\ref{sec:diff-shape}. On each pair, we ran our highly optimised,
state of the art APTED implementation, with a one second runtime budget enforced using common operating
system runtime limits.

\begin{figure*}
  \centering
  \includegraphics[width=\linewidth]{figures/time_budget}
  \caption{The one-second budget, read the same way twice. Left: the share of sampled file pairs
  whose whole-tree APTED computation completes within the time on the axis, per artifact
  category; the height of each curve at the dashed line is RA2, and the pairs the budget killed
  are the gap to 100\%. Right: wall-clock per diff for every tool configuration of
  Section~\ref{sec:tools}, every repeated run over every fixture, with \textsc{CodeDiff} charged
  for its parse; the dotted curve is the tree-sitter parse alone, the floor every AST-aware tool
  pays first. The dashed line is the same one second in both panels.}
  \label{fig:time-budget}
\end{figure*}

The left panel of Figure~\ref{fig:time-budget} shows, for each artifact category, the share of
pairs completed as a function of the time allowed. Size decides the outcome: within one second the
computation finishes for \RqOneCodeThirtyToHundredPct{}\% of code file pairs at 31 to 100 lines
and for \RqOneCodeHundredToThreeHundredPct{}\% at 101 to 300, and the curve keeps climbing at
the budget rather than flattening, so a larger budget would buy more pairs at every size. The shape
of the input matters as well: scripting languages complete more often than code, and data files
less often.

\begin{resultbox}{RA2 (Speed)}
  Within one second, pure tree edit distance algorithms find the diff for
  \RqOneCodeThirtyToHundredPct{}\% of code file pairs at 31 to 100 lines; at 101 to 300 lines it
  already falls to \RqOneCodeHundredToThreeHundredPct{}\%. Over all code pairs sampled, only
  \RqOneCodePct{}\% complete under one second.
\end{resultbox}

\section{Uniqueness}
\label{sec:uniqueness}

The Agreed and Faithful-on-screen properties fail in the same way: the tool shows one answer where
more than one is correct, and the reader has no way to know. This section measures how often that
happens, first in the mapping and then in its rendering.
Existing tools all provide a single diff. Line based tools mark lines as inserted, deleted or match lines
1:1. Tree based tools and algorithms have 4 basic operations: insert node, delete node, update node
and match nodes 1:1. Some algorithms extend the basic operations with tree equivalents, e.g. insert tree or
delete tree. No tool allows for multiple solutions.

Manually examining real-world inputs, we found that some diffs have multiple equally correct
solutions: multiple equally correct 1:1 mappings within a given many-to-many node mapping must
cover one side of the mapping completely, and the remaining nodes on the other side should be
marked as inserted or deleted, as appropriate.

Figure~\ref{fig:ambiguous-example} is a motivating example.
\begin{figure}
  \small
  \begin{tabular}{@{}p{0.46\linewidth}@{\hspace{0.04\linewidth}}p{0.46\linewidth}@{}}
    \textsc{Before} & \textsc{After} \\
    \midrule
    \begin{minipage}[t]{\linewidth}
\begin{verbatim}
def is_even(x):
  return x > 0 \
    and x % 2 == 0
\end{verbatim}
    \end{minipage}
    &
    \begin{minipage}[t]{\linewidth}
\begin{verbatim}
def is_even(x):
  return isinstance(x, int) \
    and x > 0 \
    and x % 2 == 0
\end{verbatim}
    \end{minipage} \\
  \end{tabular}
  \caption{A change with no unique correct mapping. The before AST holds two nested
  \texttt{binary\_expression} nodes and the after AST three; every choice of which two to match
  costs the same and renders the same.}
  \label{fig:ambiguous-example}
\end{figure}
The AST of the before snippet contains two \texttt{binary\_expression} nodes, one a direct child of
the other. The after snippet contains three. There is no objective way to determine which of the
three is added and which two should be matched. The edit distance cost of all three mappings is
exactly equal, and since \texttt{binary\_expression} is an internal node with no rendering of its
own, all three also produce the same visualisation---so human judgement cannot break the tie either.

\textbf{RQ3.1.} What percentage of real-world code changes have multiple optimal AST node
mappings?

We count a change as admitting several optimal mappings when its ground truth records at least one
\emph{multi-map group}: a set of nodes on one side and a set on the other that the annotator judged
interchangeable, so that any consistent pairing within the group is equally correct. Across the corpus the annotators recorded \AmbiguityGroups{} such groups.

Even without such groups, multiple solutions can exist. Where the annotator's mapping and the
matcher's mapping for one change cost exactly the same under the unit cost model and still disagree
about which nodes pair with which, the cost function cannot separate two different answers. That happens on \CostTieDifferentFixtures{} changes (\CostTieDifferentPct{}\%).

\begin{resultbox}{RA3.1 (Mapping uniqueness)}
  In at least \AmbiguityListPct{}\% of changes (\AmbiguityListWith{} of \AmbiguityListScored{}), humans found that some set of nodes admits several equally correct pairings. Independently of that annotation, on
  \CostTieDifferentPct{}\% of changes the human and algorithmic mappings tie on cost while
  disagreeing on the mapping, so the cost function alone cannot pick between them.
\end{resultbox}

\textbf{RQ3.2.} What percentage of real-world code changes have multiple equally correct textual
visualisations of the diff?

Settling the mapping does not settle what the reader sees. Of the \PaintingFixtures{} painted
changes, \PaintingDual{}---\PaintingDualPct{}\%---carry at least two paintings, the annotator's
finding that there are multiple correct renderings of the same change.

\begin{resultbox}{RA3.2 (Visualisation uniqueness)}
Of the \PaintingFixtures{} painted changes taken from real commits, \PaintingDual{} (\PaintingDualPct{}\%) admit at least two equally correct renderings.
\end{resultbox}



\subsection{How much of the mapping reaches the reader}

\textbf{RQ3.3.} What percentage of the AST influences the textual visualisation?

A syntax tree contains two kinds of nodes. Some carry text of their own---an identifier, an
operator, a literal, a comment. Others are pure structure: a \texttt{block} whose extent is
exactly its children's, an \texttt{expression\_statement} wrapping one expression, the chain of
single-child nodes a grammar interposes between a statement and the name inside it. A rendered diff highlights regions of source text, so a node of the second kind does not appear in the output on its own. A tool that maps every displayed token correctly and every structural wrapper wrongly produces a diff that is still correct; the same tool scores badly on a metric that counts all nodes equally. A node is visible if any byte
inside its extent lies outside every one of its children---that is, if it contributes text of its own to the rendering.

\begin{resultbox}{RA3.3 (Visible share of the AST)}
\InvisibleNodeShare{}\% of the AST nodes in \NumFixtures{} fixtures carry no text of their own and cannot appear in a rendered diff at all: \VisibleNodeShare{}\% are visible. Roughly one AST node in three is scaffolding.
\end{resultbox}

The consequence is not that node-level metrics should be abandoned. It is that a rate over all
nodes and a rate over visible nodes answer different questions, and any paper reporting one should
say which. This paper reports both wherever it reports an accuracy figure.

\section{State of the Art}
\label{sec:tools}

The sections above measure the problem; this one asks whether the tools that exist have the four
properties of Section~\ref{sec:intro}. We
run every established diffing tool we could obtain over the corpus of Section~\ref{sec:dataset},
score each against the same human mapping, and time each on the same machine, so that accuracy
and cost can be read together rather than one tool at a time. Two questions follow.

%% RQ4.1 and RQ4.2 are stated here, next to the measurements that answer them. Section 1 names
%% them by number under the property each one tests; keep the numbers in step.
\textbf{RQ4.1.} How accurately do current state-of-the-art code diffing tools find optimal
visualisations?

\textbf{RQ4.2.} How fast are real-world state-of-the-art tools?

\textbf{Comparison tools.} We measure seven established tools, all introduced in
Section~\ref{sec:background}, on the \NumFixtures{}-fixture corpus described in
Section~\ref{sec:dataset}. \texttt{git diff} is measured
under each of its four algorithms, Myers, minimal, patience and histogram, giving ten configurations in all. Five configurations report at line granularity only: Unix
\texttt{diff} and the four \texttt{git} algorithms. Five report sub-line detail: BDiff~\cite{bdiff},
Neovim's diff mode, and the three AST-aware tools, GumTree (\GumTreeVersion{})~\cite{gumtreediff-2014},
difftastic~\cite{difftastic} and diffsitter~\cite{diffsitter}. Our own tool's figures are reported in Section~\ref{sec:codediff}, on the same basis and against the same ground truth.



\textbf{Line-level agreement.} None of the ten configurations reports an AST-node mapping in a
form the ground truth can be compared with directly, so we compare them on a common ground: which
source lines each marks as touched, against the lines the human mapping touches. Coverage varies
widely: of the \UnixDiffFixtures{} fixtures this comparison was run over, the text-based tools
cover all of them, while diffsitter's compiled-in grammar set matches \DiffsitterFixtures{},
GumTree's generator set \GumTreeFixtures{}, and difftastic's \DifftasticFixtures{}. Because those
subsets differ so much, percentages are used.

We report agreement per fixture rather than pooled over lines in
Table~\ref{tab:agreement-buckets}, and Figure~\ref{fig:agreement-by-dataset} splits the same
reading by the dataset each fixture was sampled from.
\emph{Perfect} means literally zero mismatched lines. A diff a reader can trust on sight is
categorically different from one that is merely close. $\ge$99\% is ``a reader has to look
twice''; 95--99\% is ``a reader has to check''. Below 95\% the tool is no longer usable.

\begin{resultbox}{RA4.1 (Tool accuracy)}
No tool maps much more than three fifths of the corpus perfectly. The choice of line-diff
algorithm barely registers: Unix \texttt{diff} and all four \texttt{git} algorithms sit within a
percentage point of one another. Syntax-awareness alone does not guarantee that a tool recovers a
change the way a human reads it: diffsitter, at \DiffsitterPerfectPct{}\%, is closer to the
line-based tools than to the other syntax-aware tools.
\end{resultbox}

../../../plots/benchmark_other_buckets_combined.tex

\begin{figure*}
  \centering
  \includegraphics[width=\linewidth]{figures/benchmark_other_buckets_by_dataset}
  \caption{The line-level agreement of Table~\ref{tab:agreement-buckets}, split by the dataset
  of Section~\ref{sec:dataset} each fixture was drawn from, with the pooled reading on the right.
  Each bar is one configuration over the fixtures it supports ($n$ beside the bar); the segments
  are the four agreement buckets, and the Perfect segment carries its share.}
  \label{fig:agreement-by-dataset}
\end{figure*}

The split by dataset in Figure~\ref{fig:agreement-by-dataset} orders the datasets the same way for
every configuration: each scores its own best on either the Stratified sample or the Defects4J
changes, and all five line-granularity configurations do best on Defects4J. Bug fixes are short,
localised edits, which is the shape a line-based diff recovers most reliably, so the ordering is a
property of the datasets rather than of any tool.

\textbf{Cost.} Table~\ref{tab:speed-vs-others} reports wall-clock percentiles, sorted fastest to slowest by median. Each tool is measured over
the fixtures it supports rather than over one common subset, so the populations behind these rows
differ---the text-based tools draw on every fixture, while an AST tool draws only on the languages
its grammars cover. Two tools are reported twice: GumTree once invoked as a fresh process per fixture, the realistic cost of a single CLI
invocation, and once run inside one persistent JVM across all fixtures, isolating JVM startup from
GumTree's own matching cost; and BDiff likewise per process and inside one persistent Python
interpreter, for the same reason.

\begin{table}
  \caption{Wall-clock time percentiles for the ten tool configurations, milliseconds, sorted
  fastest to slowest by median, with \textsc{CodeDiff} on the last row. \emph{Max} is the slowest
  single one of the \TimingRepeats{} runs over any fixture. Every configuration is timed as a whole
  process except \textsc{CodeDiff}, which is timed in-process and so pays no process startup; see
  Section~\ref{sec:threats}.}
  \label{tab:speed-vs-others}
  \small
  \begin{tabular}{lrrrr}
    \toprule
    Tool & p50 & p90 & p99 & Max \\
    \midrule
    Unix \texttt{diff} & \SpeedUnixDiffPFifty{} & \SpeedUnixDiffPNinety{} & \SpeedUnixDiffPNinetyNine{} & \SpeedUnixDiffMax{} \\
    \texttt{git} (minimal) & \SpeedGitMinimalPFifty{} & \SpeedGitMinimalPNinety{} & \SpeedGitMinimalPNinetyNine{} & \SpeedGitMinimalMax{} \\
    \texttt{git} (patience) & \SpeedGitPatiencePFifty{} & \SpeedGitPatiencePNinety{} & \SpeedGitPatiencePNinetyNine{} & \SpeedGitPatienceMax{} \\
    \texttt{git} (histogram) & \SpeedGitHistogramPFifty{} & \SpeedGitHistogramPNinety{} & \SpeedGitHistogramPNinetyNine{} & \SpeedGitHistogramMax{} \\
    \texttt{git} (Myers) & \SpeedGitMyersPFifty{} & \SpeedGitMyersPNinety{} & \SpeedGitMyersPNinetyNine{} & \SpeedGitMyersMax{} \\
    BDiff (warm interpreter) & \SpeedBDiffWarmPFifty{} & \SpeedBDiffWarmPNinety{} & \SpeedBDiffWarmPNinetyNine{} & \SpeedBDiffWarmMax{} \\
    diffsitter & \SpeedDiffsitterPFifty{} & \SpeedDiffsitterPNinety{} & \SpeedDiffsitterPNinetyNine{} & \SpeedDiffsitterMax{} \\
    GumTree (warm JVM) & \SpeedGumTreeWarmPFifty{} & \SpeedGumTreeWarmPNinety{} & \SpeedGumTreeWarmPNinetyNine{} & \SpeedGumTreeWarmMax{} \\
    difftastic & \SpeedDifftasticPFifty{} & \SpeedDifftasticPNinety{} & \SpeedDifftasticPNinetyNine{} & \SpeedDifftasticMax{} \\
    \texttt{nvim -d} & \SpeedNvimDiffPFifty{} & \SpeedNvimDiffPNinety{} & \SpeedNvimDiffPNinetyNine{} & \SpeedNvimDiffMax{} \\
    BDiff (cold, per-invocation) & \SpeedBDiffColdPFifty{} & \SpeedBDiffColdPNinety{} & \SpeedBDiffColdPNinetyNine{} & \SpeedBDiffColdMax{} \\
    GumTree (cold, per-invocation) & \SpeedGumTreeColdPFifty{} & \SpeedGumTreeColdPNinety{} & \SpeedGumTreeColdPNinetyNine{} & \SpeedGumTreeColdMax{} \\
    \midrule
    \textsc{CodeDiff} & \SpeedCodeDiffPFifty{} & \SpeedCodeDiffPNinety{} & \SpeedCodeDiffPNinetyNine{} & \SpeedCodeDiffMax{} \\
    \bottomrule
  \end{tabular}
\end{table}

\begin{resultbox}{RA4.2 (Tool cost)}
The ten configurations span more than two orders of magnitude at the median, from Unix
\texttt{diff}'s \SpeedUnixDiffPFifty{}\,ms to GumTree's \SpeedGumTreeColdPFifty{}\,ms per
invocation. At the 99th percentile the same range runs from \SpeedUnixDiffPNinetyNine{}\,ms to
\SpeedGumTreeColdPNinetyNine{}\,ms, three orders of magnitude apart.
\end{resultbox}

We report both a cold and a warm figure for GumTree and BDiff because they answer different questions: the cold number is what a developer running a diff from the command line waits for, and the warm number is the matching cost once process startup is amortised away.

Every timing number in this paper is the result of repeated runs rather than a single
measurement. Each tool is run \TimingRepeats{} times on every fixture, and the percentiles are
taken over every individual run pooled across fixtures rather than over per-fixture means: the
repeat count is uniform, so no fixture is over-weighted, and the spread within a fixture survives
into the number instead of being averaged away before it is ever read.



\section{\textsc{CodeDiff}}
\label{sec:codediff}

This section is the paper's tool contribution. \textsc{CodeDiff} is not a new matching
algorithm: every technique in it is published, and Section~\ref{sec:related} says where. What is
new is the order in which those techniques run and what each is allowed to decide, and every step
of that order is a consequence of a measurement above. Table~\ref{tab:findings-to-design} lists
the correspondence; the phases that follow are the same list read top to bottom, and the tool is
then evaluated against the same ground truth and on the same basis as Section~\ref{sec:tools}'s
ten configurations.

\begin{table}
  \caption{What each empirical result decided in \textsc{CodeDiff}'s design.}
  \label{tab:findings-to-design}
  \small
  \begin{tabular}{@{}>{\raggedright\arraybackslash}p{0.44\linewidth}>{\raggedright\arraybackslash}p{0.52\linewidth}@{}}
    \toprule
    Finding & Consequence \\
    \midrule
    The median edit rewrites \EditsChurnPFiftyPct{}\% of its file (Section~\ref{sec:edit-shape})
      & Phase 1 settles the unchanged majority by hashing before any edit-distance computation
        runs \\
    Whole-tree APTED finishes within one second for only \RqOneCodePct{}\% of code pairs (RA2)
      & APTED runs only inside one already-matched pair (Phase 3), never over the whole tree \\
    On \CostHumanHigherPct{}\% of changes the human mapping costs more than the cheapest one
      (RA1.2)
      & Phases 2 and 3 encode reader preferences that cost more than the optimum: comments and
        modifiers follow their node, names beat cost, diagnostics are held back \\
    \AmbiguityListPct{}\% of changes have no expressible optimum and \CostTieDifferentPct{}\%
      tie at different mappings (RA1.1, RA3.1)
      & The phase order is the tie-break: the earliest phase to claim a node wins, and the
        closing step marks what no phase claimed \\
    \InvisibleNodeShare{}\% of nodes never reach the screen (RA3.3)
      & Accuracy is reported on visible nodes and on rendered lines, not on the node count alone \\
    No established tool renders more than \GumTreePerfectPct{}\% of changes perfectly (RA4.1)
      & The bar the evaluation has to clear, on the same basis \\
    Files reach \AstMax{} nodes (Section~\ref{sec:dataset})
      & The Robust target below \\
    \bottomrule
  \end{tabular}
\end{table}

Two of those consequences are numeric targets rather than design decisions, and the paper states
them as such:

\begin{itemize}
  \item \textbf{Robust:} \textsc{CodeDiff} must handle every file up to the measured maximum,
    \AstMax{} AST nodes, with headroom above that number.
  \item \textbf{Fast:} \textsc{CodeDiff} must compute a diff in under \SpeedTargetMs{}\,ms for
    \SpeedTargetPct{}\% of real-world commits, the Timely property's one second applied to the
    commit population rather than to a corpus of hard cases.
\end{itemize}

\textsc{CodeDiff} matches two ASTs in \NumPhasesWord{} phases. Each phase either resolves part of the
mapping directly or narrows the residual left for the next phase. Later phases only ever see
nodes earlier phases left unmatched. A sixth step closes the mapping, recording the deletion or
insertion implied by whatever every phase above left undecided.

\textbf{Phase 1, hash-based descent.} \textsc{CodeDiff} matches subtrees by two hash variants, in
order: byte-identical subtrees, then same-shape subtrees with any leaf value. This phase is
motivated directly
by the commit-size data of Section~\ref{sec:dataset}: since most commits touch a small fraction of
a file's lines, most of a typical file's AST is byte-identical between before and after. Resolving
that identical majority with cheap hashing, before any tree-edit-distance computation runs at all,
avoids paying that computation's cost on content that never changed.

\textbf{Phase 2, contextual exact matching.} Comments and modifiers (attributes, decorators) that
immediately precede an already-matched node, and diagnostic statements (logging calls, assertions,
panics) that are byte-identical on both sides, get matched directly. Neither is reachable by phase 1's structural hashing: comments and modifiers are not part of the hashed AST
shape, and diagnostic statements are held back deliberately, so that matching one does not
fragment a larger match around it. These are the preferences RA1.2 measures: each of them can
cost more than the cheapest mapping, and each is what a reader marks.

\textbf{Phase 3, syntax-aware subtree matching.} This phase generalises three matching strategies:
matching by fully-resolved name (handling overloads and duplicate names across
scopes), a greedy cost-estimate matcher for containers with no name structure, such as an
\texttt{if} block or a loop body, and an overlap matcher pairing import statements that share a
base path. Before any of these run, large flat subtrees, such as a macro's token list or a big
flat argument list, are pre-matched directly with Myers' algorithm~\cite{myers-1986}, since a flat sequence has no tree structure worth exploiting. Each accepted pair is then handed to
APTED~\cite{apted-2015}, scoped to that pair alone, to compute its exact internal edit script.
This is RA2 applied: the exact computation is affordable inside one matched pair, and only there.

\textbf{Phase 4, residual alignment.} Whatever remains unmatched after phases 1 through 3 is
resolved by a Myers-LCS alignment~\cite{myers-1986} over the residual forest, bracketed by two
passes of strict bottom-up propagation: a node whose children all match the children of exactly
one node on the other side is matched to it, which shrinks the residual, and is run again afterward so that pockets the alignment just resolved can still bubble up to
their now-complete ancestors.



\textbf{Phase 5, move-detection recovery.} The final matching pass. Ordered tree edit distance cannot
express a subtree that relocated across a matched boundary as anything but a delete plus an
insert. This phase, adapted directly from GumTree's recovery-mappings
phase~\cite{gumtreediff-2014}, pairs up whole subtrees left fully deleted on one side and fully
inserted on the other, when they are byte-identical and large enough to rule out coincidence.
Move is the operation the introduction argues a diff needs and line-based tools lack, so the
mapping is not closed until this pass has run.

\textbf{Result.} \textsc{CodeDiff} renders \CodeDiffPerfectPct{}\% of the corpus perfectly,
against \GumTreePerfectPct{}\% for the best established configuration, at a median of
\SpeedCodeDiffPFifty{}\,ms including its parse, and it completes every fixture. Where it falls short
is where every tool falls short: the Full list and the largest files. The rest of this section is
the evidence for that sentence, then its limits.

\textbf{Accuracy.} Against the human mapping, \textsc{CodeDiff} maps \NodeAccuracyPct{}\% of AST
nodes correctly (\NodeMismatches{} mismatches in \NodesTotal{}), and \VisibleNodeAccuracyPct{}\% of
the \VisibleNodeShare{}\% of nodes that carry text and can reach the screen
(\VisibleNodeMismatches{} in \VisibleNodesTotal{}). On the line-level basis of
Section~\ref{sec:tools} it marks \CodeDiffLineMismatches{} lines differently from the human
mapping over \CodeDiffFixtures{} fixtures, a rate of \CodeDiffLineRate{}\%, the lowest of any tool
measured, and it moves only to \CommonCodeDiffLineRate{}\% on the \CommonFixtures{}-fixture subset
every tool scored. Table~\ref{tab:codediff-by-dataset} splits the result by dataset in
the buckets of Table~\ref{tab:agreement-buckets}: the perfect share is
\CodeDiffPerfectPctStratified{}\% on the Stratified sample and \CodeDiffPerfectPctFull{}\% on the
Full list, and the line rate (\CodeDiffLineRateStratified{}\% against \CodeDiffLineRateFull{}\%),
node rate (\CodeDiffNodeRateStratified{}\% against \CodeDiffNodeRateFull{}\%) and visible-node
rate (\CodeDiffVisibleRateStratified{}\% against \CodeDiffVisibleRateFull{}\%) order the datasets
the same way.

../../../plots/benchmark_codediff_by_dataset.tex

\textbf{Robustness.} Run over the \RobustnessFixtures{} fixtures with a
\RobustnessTimeoutSeconds{}-second budget per fixture and no cap on tree size,
\textsc{CodeDiff} completes every one: \RobustnessTimedOut{} timeouts and \RobustnessPanicked{}
panics across \RobustnessLanguages{} languages. The largest fixture holds \RobustnessMaxNodes{}
AST nodes per side, the slowest takes \RobustnessSlowestMs{}\,ms, and peak heap use is
\RobustnessPeakMemoryMb{}\,MB. The same measurement over the whole Full corpus reaches the files
the fixture corpus does not. Every modified code file in the corpus's recent history
(Section~\ref{sec:edit-shape}), \RobustnessFullAttempted{} pairs across
\RobustnessFullLanguages{} languages, was diffed with no node cap under a
\RobustnessFullTimeoutSeconds{}-second budget and a \RobustnessFullMemoryCapGb{}\,GB memory cap
per process. \textsc{CodeDiff} completes \RobustnessFullCompletedPct{}\% of them with
\RobustnessFullPanicked{} panics, at a median of \RobustnessFullPFiftyMs{}\,ms and a 99th
percentile of \RobustnessFullPNinetyNineMs{}\,ms; the largest completed pair holds
\RobustnessFullMaxSideNodes{} AST nodes on each side, against the \AstMax{} of the largest file
in the corpus that the Robust target names. The remainder is \RobustnessFullTimedOut{} pairs past
the budget and \RobustnessFullKilled{} past the memory cap. The killed pairs are commits of
\RobustnessFullKilledFiles{} distinct files, all but one of them generated parser tables,
minified bundles or embedded binary data; the exception is one commit of a 40,000-line
single-header C++ library. Given \RobustnessFullBigMemoryCapGb{}\,GB and a five-minute budget, one
commit of each memory-killed file was tried again: \RobustnessFullKilledFilesCompletedUnderBigCap{}
completed, one needed more memory still, and the rest ran past even that budget.

\textbf{Speed.} Over the \CodeDiffFixtures{} fixtures, \textsc{CodeDiff}'s median wall-clock time
is \SpeedCodeDiffPFifty{}\,ms, its 90th percentile \SpeedCodeDiffPNinety{}\,ms and its 99th
percentile \SpeedCodeDiffPNinetyNine{}\,ms, parse included. The right panel of
Figure~\ref{fig:time-budget} places that curve among the ten configurations of
Table~\ref{tab:speed-vs-others}: the line-based tools sit to its left, the AST-aware tools to its
right, and it is the only AST-aware curve that reaches 100\% before the budget line.

\textbf{Limits.} Three things the result does not show. The slowest single run of the
\TimingRepeats{}, \SpeedCodeDiffMax{}\,ms on the largest fixture in the corpus, is past the
\SpeedTargetMs{}\,ms target; the target is stated over real-world commits and the corpus was
selected for hard changes, but the miss is real, and the other configurations miss the same
fixture by one to two orders of magnitude. The Robust target is not met at its ceiling: the
corpus holds files of up to \AstMax{} nodes, the largest pair the whole-corpus run diffed to
completion has \RobustnessFullMaxSideNodes{} a side, and the larger ones it met, all generated
tables, needed more memory than the cap rather than more time (Section~\ref{sec:threats}). And
\textsc{CodeDiff} is timed
in-process while every other configuration is timed as a whole process, so its curve pays no
process start-up; Section~\ref{sec:threats} bounds that asymmetry.

\section{Threats to Validity}
\label{sec:threats}

\textbf{Construct validity.} Every rate this paper reports about ambiguity is a floor rather than
an estimate. A multi-map group exists where an annotator noticed that two pairings were equally
correct and recorded it; a change whose alternative mapping nobody saw is filed as unique. The
same holds for the paintings. We report these as existence results and lower bounds throughout, and RA3.1's second reading---cost ties at differing mappings, which needs no annotation to detect---exists precisely because it does not depend on annotator attention. For the same reason we score RA1.1 and RA3.1 over the \AmbiguityListScored{} changes whose ambiguity pass is complete rather than over all \NumFixtures{}. That pass covers the Curated and Full lists in full, so for those two lists the bound is one on annotator attention rather than one on coverage. The Stratified sample (\AmbiguityStratifiedWith{} groups in \AmbiguityStratifiedTotal{} changes) and the Defects4J units (\AmbiguityDefectsFourJWith{} in \AmbiguityDefectsFourJTotal{}) run far under the rate the two finished lists show, which is an annotation gap rather than a finding about those changes; Table~\ref{tab:expressibility} prints both rows and its \emph{All} row under that caveat.

Because no whole-file exact computation is affordable at these sizes (RA2), RA1.2 compares the
human mapping's cost to a heuristics-augmented computation's rather than to a proven optimum. Its
rate is therefore a lower bound.

Section~\ref{sec:tools} scores tools at the line level because that is the only granularity all
ten configurations share. A line-level projection of a node mapping discards detail that the
AST-aware tools compute and the line-based ones never had, which if anything favours the
line-based tools. It also means the paintings---the only place the corpus records a second correct
rendering---never enter that comparison, so RA3.2's error term sits in those tables unremoved.

\textbf{Internal validity.} \textsc{CodeDiff} is timed in-process while every external tool is
timed as a whole process, so it pays no process-startup cost that the others pay in full. We state
this rather than correct for it, and report GumTree and BDiff both cold and warm so that a reader
can see how large the effect is for the two tools where it dominates. Per-tool mean runtimes are
additionally distorted by run order, whichever tool runs first per fixture absorbing cold-cache
cost; to minimise this effect we quote percentiles rather than means throughout, and every timing number is the result of repeated runs.

The ground truth is authored by the same people who built \textsc{CodeDiff}, which is a real risk
of bias in its favour. Three things limit it: the mappings are authored against the change itself in
a dedicated annotation tool with no tool output shown, the full annotation history is published
with the corpus so that any individual mapping can be inspected and disputed, and
Section~\ref{sec:oracle} reads our annotation, and the tool over the same units, against ground
truth authored by other people, for a different tool, before this one existed. That last check is the only one of the three that does not
depend on our own discipline, and it is the reason the corpus was extended with Defects4J at all.

\textbf{External validity.} The fixture corpus is deliberately not a fully random sample of
real commits. It concentrates hard, structurally interesting changes, because easy ones exercise
nothing worth measuring---so every tail figure in this paper is a stress reading rather than an
estimate of what a developer waits for on a typical commit. Every accuracy result is also reported per dataset (Figure~\ref{fig:agreement-by-dataset},
Table~\ref{tab:codediff-by-dataset}), so a reader can separate the two random samples, the
stratified one, and the Defects4J bug fixes from each other. The Defects4J rows carry the
selection caveat of Section~\ref{sec:diff-shape}: the units solved so far are how far annotation
has reached, not a draw, and they run shorter than the ones still unsolved.

The painted subset is smaller still and is not randomly drawn: painting is manual, and which
changes have been painted is an artifact of the order the annotators worked in. It runs to
shorter changes than the corpus as a whole, so RA3.2's rate may move in either direction as
painting reaches larger changes.

Finally, the fixture corpus's largest file holds \RobustnessMaxNodes{} AST nodes, well short of
the Robust target's \AstMax{}. The target is exercised only by the whole-corpus run of
Section~\ref{sec:codediff}, whose budget and memory cap were applied with eight processes
sharing the machine, so a pair near either limit may have crossed it for want of a core rather
than of an algorithm. Re-running the timeouts on a quiet machine moved two of them into the
completed set and no more, and the memory kills recur across every commit of the same files,
so neither limit looks like an artefact of the sharing. The empirical-study
numbers of Section~\ref{sec:dataset} do cover the full \NumRepos{}-repository list, but the
\NumFixtures{} solved changes are drawn from it by the sampling of
Section~\ref{sec:diff-shape} rather than being representative of it by construction.

\section{Related Work}
\label{sec:related}

\textbf{Tree-diffing tools.} Section~\ref{sec:background} describes what each of the tools
compared here does; what matters for positioning this work is where each one's goal differs from
ours. GumTree~\cite{gumtreediff-2014} remains the reference point, and \textsc{CodeDiff}'s
move-recovery phase adapts its recovery-mappings heuristic directly: the difference is that GumTree
targets research use, with broad language and backend flexibility that stays valuable well beyond
what RA4.1's single corpus measures, while \textsc{CodeDiff} narrows to a tree-sitter-only frontend
and spends the saved generality on production targets. difftastic~\cite{difftastic} and
diffsitter~\cite{diffsitter} share that parsing layer but optimise for a diff a human reads
comfortably rather than for fidelity of the underlying node mapping, so RA4.1 measures them against
an objective they did not choose. Against all of them, Unix \texttt{diff} and \texttt{git
diff}~\cite{myers-1986} remain the baseline: on our corpus they are the fastest measured and, at
the line level, more accurate than every pre-existing AST-aware tool compared. BDiff's~\cite{bdiff}
block awareness does not move that metric, and neither does the choice among \texttt{git}'s four
algorithms~\cite{nugroho-histogram}.

\textbf{Ground truth and its ambiguity.} GumTree's own evaluation, and most that followed it,
score a tool against a single reference mapping per change. RA1.1, RA3.1 and RA3.2 are an argument
that this default understates what a corpus can be asked: on the changes where the correspondence
is not one-to-one, where several mappings are equally correct, or where several renderings are, a
fixed reference records a correct answer as wrong. The corpus here departs from the default in two
ways---a multi-map group, which records an $N{:}M$ site and admits any consistent one-to-one
pairing of it, and a second, independently authored painting where the rendering is not
unique---and both were added because annotating real changes forced them, not by design. We are
not aware of a code-diffing ground truth that records either. RA1.2 adds a third departure from
the usual assumptions: the unit cost model that ranks candidate mappings disagrees with the
annotators on a measurable minority of changes.

\textbf{Empirical studies.} Dyer et al.~\cite{ast-api-1} and L\"{a}mmel et al.~\cite{ast-api-2}
mine AST nodes at scale to study API usage; Arafat and
Riehle~\cite{commit-size-distribution-paper} measure the commit-size distribution of open-source
software, the distribution \textsc{CodeDiff}'s Fast target is stated against; and
Defects4J~\cite{defects4j} collects real, reproducible bug fixes from open-source Java projects,
each a before/after file pair, and has since served as the file-pair corpus behind most
evaluations of AST differencing on Java. This paper's empirical study differs from these in
purpose rather than method: it measures file- and AST-size
distributions specifically to set, and then evaluate, the engineering targets of a diffing tool.

\section{Conclusions}
\label{sec:conclusion}

Line-based and AST-based diffing tools can meet developers' speed requirements: the line-based
tools answer within tens of milliseconds at every percentile measured, and the fastest AST-aware
tools stay inside the one-second bound across the corpus (RA4.2). Current algorithms and implementations fall short of meeting the
full quality requirements that real changes impose. On \AmbiguityListPct{}\% of changes the
best possible correspondence cannot be expressed with the four standard operations at all
(RA1.1), on \CostHumanHigherPct{}\% the cheapest mapping is not the one humans judge correct
(RA1.2), the mapping and its rendering are not unique (RA3.1, RA3.2), and no established tool
maps much more than half of the corpus perfectly (RA4.1). Closing that gap needs a ground truth
that records the correspondence itself, a cost model that agrees with readers, and tools that
are measured against both; the corpus published with this paper and the \textsc{CodeDiff} tool
are a first step towards all three.

\bibliographystyle{ACM-Reference-Format}
\bibliography{references}

\end{document}
